% ======================================================================
%  ANNEXES -- slides de secours, à montrer seulement si une question
%  porte dessus (non comptées dans la numérotation principale)
% ======================================================================

% ----------------------------------------------------------------------
% SI QUESTION sur le réglage de la loop closure
%  - Le réglage de ces paramètres, surtout le rayon de recherche et le
%    seuil de fitness, fixe la sensibilité de la détection : rayon trop
%    grand ou seuil trop permissif -> fausses détections ; valeurs trop
%    strictes -> vraies boucles manquées.
%  - Sur KITTI (bag sans vraie boucle), j'ai forcé une loop closure en
%    réduisant fortement le délai minimal (historyKeyframeSearchTimeDiff),
%    uniquement pour tester la communication framework <-> LIO.
% ----------------------------------------------------------------------
\begin{frame}{Annexe : paramètres de loop closure de LIO}
  \centering
  \small
  \begin{tabular}{@{}l c l@{}}
    \toprule
    \textbf{Paramètre} & \textbf{Défaut} & \textbf{Rôle} \\
    \midrule
    \texttt{loopClosureEnableFlag}         & --     & active / désactive la détection \\
    \texttt{loopClosureFrequency}          & 1 Hz   & fréquence des tentatives \\
    \texttt{historyKeyframeSearchRadius}   & 15 m   & rayon de recherche des candidates \\
    \texttt{historyKeyframeSearchTimeDiff} & 30 s   & écart temporel minimal \\
    \texttt{historyKeyframeSearchNum}      & 25     & keyframes par sous-carte (ICP) \\
    \texttt{historyKeyframeFitnessScore}   & 0,3    & seuil ICP (plus bas = plus strict) \\
    \bottomrule
  \end{tabular}\par
  \vspace{0.4cm}
  {\small rayon trop grand / seuil permissif $\rightarrow$ \alert{fausses boucles}\\[2pt]
   valeurs trop strictes $\rightarrow$ \alert{boucles manquées}\par}
\end{frame}

% ----------------------------------------------------------------------
% SI QUESTION sur le champ ring (KITTI)
%  - Chaque nappe du LiDAR correspond à un angle d'élévation fixe ; on
%    rattache le point à la nappe dont l'angle est le plus proche.
%  - D = distance du point au LiDAR, theta = arcsin(z / D), puis on ramène
%    theta dans [0, N] entre theta_min = -24,8° et theta_max = 2°
%    (Velodyne de KITTI, N = 64).
%  - Hypothèse : espacement angulaire uniforme entre nappes -> c'est une
%    approximation.
% ----------------------------------------------------------------------
\begin{frame}{Annexe : calcul du champ \texttt{ring} (KITTI)}
  \begin{columns}[c]
    \begin{column}{0.52\textwidth}
      \centering
      \begin{tikzpicture}[scale=1.05, >={Stealth}]
        \coordinate (O) at (0,0);
        \foreach \ang in {-15,-9,-3,3,9,15,21} {
          \draw[gray!55, thin] (O) -- (\ang:4.6);
        }
        \node[gray!70] at (-15:5.05) {\scriptsize ring $0$};
        \node[gray!70] at (21:5.15) {\scriptsize ring $N{-}1$};
        \draw[dashed] (O) -- (4.8,0);
        \node[below right] at (4.8,0) {\scriptsize $\theta = 0$};
        \coordinate (P) at (6:4);
        \draw[thick, ->, accent] (O) -- (P);
        \fill[accent] (P) circle (1.8pt);
        \node[above right=1pt] at (P) {$P(x,y,z)$};
        \node[above=1pt] at (6:1.9) {\scriptsize $D$};
        \draw (0.9,0) arc (0:6:0.9);
        \node at (1.45,0.11) {\scriptsize $\theta$};
        \fill (O) circle (2pt);
        \node[left=2pt] at (O) {\scriptsize LiDAR};
      \end{tikzpicture}
    \end{column}
    \begin{column}{0.44\textwidth}
      \[
        D = \sqrt{x^2 + y^2 + z^2}
      \]
      \[
        \theta = \arcsin\!\left(\frac{z}{D}\right)
      \]
      \[
        \text{ring} = \frac{\theta - \theta_{\min}}{\theta_{\max} - \theta_{\min}} \times N
      \]
      \vspace{0.1cm}
      \begin{itemize}
        \item Velodyne KITTI : $N = 64$,\\ $\theta_{\min} = -24{,}8^\circ$, $\theta_{\max} = 2^\circ$
        \item \alert{Approximation} : nappes régulièrement espacées
      \end{itemize}
    \end{column}
  \end{columns}
\end{frame}

% ----------------------------------------------------------------------
% SI QUESTION sur le format de sortie (carte a priori)
%  - Une "feature" GeoJSON par case : centre de la case, label, identifiant
%    OSM d'origine ; le polygone de la case reste disponible.
%  - Système de coordonnées : Lambert-93 (EPSG:2154), ramené dans un
%    repère local en soustrayant l'origine de la carte.
% ----------------------------------------------------------------------
\begin{frame}[fragile]{Annexe : format GeoJSON (\texttt{export\_grid.geojson})}
  \begin{columns}[c]
    \begin{column}{0.66\textwidth}
\begin{semiverbatim}\scriptsize
\{
  "type": "FeatureCollection",
  "name": "export_grid",
  "crs": \{ "type": "name", "properties":
    \{ "name": "urn:ogc:def:crs:EPSG::2154" \} \},
  "features": [
    \{ "type": "Feature",
      "properties": \{
        "center_x": 0.5, "center_y": 159.5,
        "label": "road", "id": "way/154556229" \},
      "geometry": \{ "type": "Polygon",
        "coordinates": [ [ [ 0.0, 159.0 ], [ 1.0, 159.0 ],
          [ 1.0, 160.0 ], [ 0.0, 160.0 ],
          [ 0.0, 159.0 ] ] ] \} \},
    ...
  ]
\}
\end{semiverbatim}
    \end{column}
    \begin{column}{0.3\textwidth}
      \begin{itemize}
        \item 1 case = 1 \alert{feature}
        \item centre, \alert{label}, id OSM
        \item polygone de la case
        \item Lambert-93 ({\NoAutoSpacing\alert{EPSG:2154}})
      \end{itemize}
    \end{column}
  \end{columns}
\end{frame}

% ----------------------------------------------------------------------
% SI QUESTION sur la forte dérive du Scout
%  - Essai complémentaire : lecture du bag du Scout à partir de 60 s, pour
%    éviter la phase de forte dérive observée au début de l'acquisition.
%  - (À adapter : ces figures ne sont pas commentées dans le rapport, dire
%    ce que tu en as conclu.)
% ----------------------------------------------------------------------
\begin{frame}{Annexe : Scout, bag lu à partir de 60 s}
  \centering
  \begin{tabular}{@{}c@{\hspace{0.5cm}}c@{}}
    \adjincludegraphics[height=4.6cm, trim={{0.485\width} {0.26\height} {0.138\width} {0.479\height}}, clip]{scoot_loop_differe.png} &
    \adjincludegraphics[height=4.6cm, trim={{0.25\width} {0.2\height} {0.1\width} {0.12\height}}, clip]{scoot_liosam_differe_loop.png} \\
    \legende{Carte sémantique, loop closure activée} &
    \legende{Trajectoire LIO} \\
  \end{tabular}
\end{frame}
